\documentclass[10pt]{article}
%\usepackage[T1]{fontenc}
\usepackage[utf8]{inputenc}
%\usepackage[francais]{babel}
%\usepackage{listingsutf8}

\usepackage[french]{babel}
%\usepackage[T1]{fontenc}
\usepackage{comment,fullpage}
\usepackage{color}
\usepackage{graphicx}
\usepackage{lastpage}
\usepackage{amsmath}

\setlength{\parindent}{0pt}
\setlength{\parskip}{3pt}

\newcommand{\red}[1]{\textcolor{red}{#1}}

%\newenvironment{solution}{\paragraph{Correction}}{}
\newenvironment{solution}{\vspace{6pt}\noindent\color{red}\textbf{Correction : }}{\vspace{6pt}\ignorespacesafterend}

\includecomment{comment}
\excludecomment{solution}

% MATH -----------------------------------------------------------
\newcommand{\norm}[1]{\left\Vert#1\right\Vert}
\newcommand{\abs}[1]{\left\vert#1\right\vert}
\newcommand{\set}[1]{\left\{#1\right\}}
\newcommand{\Real}{\mathbb R}
\newcommand{\eps}{\varepsilon}
\newcommand{\To}{\longrightarrow}
\newcommand{\BX}{\mathbf{B}(X)}
\newcommand{\A}{\mathcal{A}}
% ----------------------------------------------------------------

\begin{document}

\begin{center}
{\Large NFP 107 -- Systèmes de gestion de bases de données}\\\vspace{1ex}
{\Large Examen juin 2018. Durée: 2h30.}\\

\end{center}

\vspace{1ex}
\hrule
\vspace{1ex}
\textbf{Consignes :}\\
-- Tous les documents sont autorisés.\\
-- Les téléphones mobiles et autres équipements communicants (exemple: PC, tablette, etc.)  doivent être éteints et rangés dans les sacs pendant toute la durée de l'épreuve.\\
-- Pensez à reporter votre numéro d'anonymat sur toutes les feuilles.\\
-- Merci de soigner votre écriture et de ne pas rédiger au crayon papier sur votre copie.\\
Ce sujet comporte 3 pages d'énoncé.

\vspace{1ex}
\hrule
\vspace{1ex}

%\input{./titlepage.tex}
% \addtocounter{page}{1}

\vspace{1ex}

Le schéma de base de données suivant sera utilisé pour l'ensemble du
sujet. Il permet de gérer une plateforme de vente de billets de spectacles en ligne :

\begin{verbatim}
SALLE(id, nom, adresse, nbPlaces)
SPECTACLE(id, idArtiste, idSalle, dateSpectacle, nbBilletsVendus)
ARTISTE(id, nom, prénom, age)
BILLET(id, idSpectateur, idSpectacle, catégorie, prix)
SPECTATEUR(id, nom, prénom, age, adresse) 
\end{verbatim}

Les artistes sont soit des individus, soit des groupes; dans le cas de groupes, le nom du groupe se trouve dans l'attribut \texttt{nom}, l'attribut \texttt{prenom} est absent et l'âge du groupe est dans l'attribut \texttt{age}.

\section{Algèbre relationnelle et SQL (9 points)}

\begin{enumerate}

\item (SQL et algèbre) Noms et prénoms des spectateurs qui ont acheté au moins une fois un billet de plus de 500 euros. (1 point)

\begin{solution}\\
Algèbre:\\
    $\pi_{nom, prenom}(SPECTATEUR \underset{id=idSpectateur}{\bowtie} \sigma_{prix > 500}(BILLET))$
\vspace{10pt}\\
SQL:
\begin{verbatim}
    SELECT DISTINCT s.nom, s.prenom
    FROM SPECTATEUR s, BILLET b
    WHERE s.id = b.idSpectateur AND b.prix > 500;
\end{verbatim}
\end{solution}

\item (SQL et algèbre) Noms des artistes qui ont donné au moins un spectacle dans la salle 'Bercy' après le 01/01/2010. (2 points)

\begin{solution}\\
Algèbre:\\
	$\pi_{ARTISTE.nom}((ARTISTE \underset{ARTISTE.id=SPECTACLE.idArtiste}{\bowtie} \sigma_{nom = 'Bercy'}(SALLE))\\ \underset{SALLE.id=SPECTACLE.idSalle}{\bowtie} \sigma_{SPECTACLE.dateSpectacle > '01/01/2010'}(SPECTACLE))$
\vspace{10pt}\\
SQL:
   \begin{verbatim}
    SELECT DISTINCT a.nom
    FROM ARTISTE a, SPECTACLE s, SALLE l
    WHERE a.id = s.idArtiste AND s.idSalle = l.id AND
          l.nom = 'Bercy' AND s.dateSpectacle > '01/01/2010';
\end{verbatim}
\end{solution}

\item (SQL seulement) Spectacle(s) pour le(s)quel(s) le nombre de billets vendus est le plus élevé. (1 point)

\begin{solution}\\
	SQL:
    \begin{verbatim}
    SELECT id
    FROM SPECTACLE
    WHERE nbBilletsVendus >= ALL (SELECT nbBilletsVendus
                                  FROM SPECTACLE);
    \end{verbatim}
\end{solution}

\item (SQL seulement) Noms des artistes et nombre total de billets vendus triés en ordre décroissant de nombre total du billets vendus. (2 points)

\begin{solution}\\
SQL:
\begin{verbatim}
    SELECT a.id, a.nom, SUM(s.nbBilletsVendus)
    FROM ARTISTE a, SPECTACLE s
    WHERE a.id = s.idArtiste
    GROUP BY a.id
    ORDER BY SUM(s.nbBilletsVendus) DESC;
\end{verbatim}
\end{solution}

\item (SQL seulement) Noms des artistes et rapport entre le nombre total de billets vendus et la somme des capacités de salles, pour les artistes pour lesquels ce rapport est inférieur à 0.8. (2 points)

\begin{solution}\\
SQL:
\begin{verbatim}
    SELECT a.id, a.nom, SUM(s.nbBilletsVendus)/SUM(l.nbPlaces)
    FROM ARTISTE a, SPECTACLE s, SALLE l
    WHERE a.id = s.idArtiste AND s.idSalle = l.id
    GROUP BY a.id
    HAVING SUM(s.nbBilletsVendus)/SUM(l.nbPlaces) < 0.8;
\end{verbatim}
\end{solution}

\item (Algèbre seulement) Noms et prénoms des spectateurs qui ont acheté au moins un billet à chaque spectacle. (1 point)

\begin{solution}\\
Algèbre:\\
    $R1 = \pi_{idSpectateur, idSpectacle} (BILLET)$\\
    $R2 = R1 \div \pi_{idSpectacle}(BILLET)$\\
    $Resultat = \pi_{nom, prenom}(SPECTATEUR \underset{SPECTATEUR.id=R2.idSpectateur}{\bowtie} R2)$
\end{solution}


\end{enumerate}

\section{Indexation (4 points)}

Dans ce qui suit, on ne considère que la variante de l'arbre B dite "B+"  dans laquelle
l'index est séparé du fichier de données.

\begin{enumerate}
  \item Combien peut-on créer d'arbres B sur une même table ? (1 pt)
 
\begin{solution}
Autant que l'on souhaite.
\end{solution}

  \item Décrire précisément la structure et le contenu de l'ensemble des feuilles d'un arbre B. Que
     contiennent-elles, dans quel ordre, comment sont-elles organisées les unes par
     rapport aux autres. (1,5 points)

\begin{solution}
      Cf. le cours. Les feuilles contiennent des entrées de la forme $(cle, adresse)$, triées sur la clé.
      Toutes les valeurs de clés sont représentées.
      Les feuilles sont chaînées les unes aux autres pour former une structure linéaire ordonnées sur les clés.

\end{solution}

\item  Un arbre B est créé sur la clé primaire qui est auto-incrémentée. On sait donc
   que la clé de chaque nouvel enregistrement est supérieure à toutes les autres.
   
   Quel impact cela a-t-il sur l'arbre B? Où se font les insertions, les éclatements (1,5 point)

\begin{solution}
Toute insertion a toujours lieu dans la feuille la plus à droite, et c'est donc toujours celle-là qui fait
l'objet d'un éclatement.
\end{solution}

\end{enumerate}


\section{Optimisation (3 points)}

On suppose que les clés primaires sont indexées, mais pas les clés étrangères.

\begin{enumerate}
\item L’optimisation de requêtes, c’est 
\begin{enumerate}
    \item Modifier une requête SQL pour qu’elle soit la plus efficace possible.
    \item Structurer les données pour qu’elles soient adaptées aux requêtes soumises.
    \item Choisir, pour chaque requête, la meilleure manière de l’exécuter.
\end{enumerate}

Justifiez brièvement votre réponse (1 point)

\begin{solution}
Réponse 3: on laisse le système choisir la meilleure manière de l’exécuter.
\end{solution}

\item Pour la première requête SQL de cet examen, indiquez quels sont les index exploitables par 
un algorithme de jointure. (1 point)

\begin{solution}
Réponse: l'index sur l'id du spectateur
\end{solution}

\item Comment expliqueriez-vous le constat que, dans la plupart des cas,  il existe un index 
exploitable pour calculer une jointure  (1 point)

\begin{solution}
Réponse: la plupart des jointures se font par égalité entre une clé primaire et une clé étrangère,
et la clé primaire est toujours indexée.
\end{solution}

\end{enumerate}

\section{Concurrence (4 points)}

On prend le schéma donné au début de l'examen.
      
\begin{enumerate}

\item Exprimez au moins un critère de cohérence de cette base (1 point)
 

\begin{solution}
Réponse: le nombre de billets vendus (\texttt{nbBilletsVendus} dans \texttt{Spectacle}) pour un spectacle est égal au nombre
de lignes dans \texttt{Billet} pour ce même spectacle.

Autre possibilité\,: pas plus de billets vendus que de places dans la salle.
\end{solution}

\item On écrit une procédure de réservation d'un billet pour un spectacle.

\begin{verbatim}
Reserver (idSpectateur, idSpectacle, categorie, prix)
   /// Requête SQL de mise à jour du spectacle
   requete A
   // Requête SQL d'insertion du billet
   requete  B
\end{verbatim}

Donnez les deux requêtes SQL, indiquez l'emplacement des \texttt{commit} et \texttt{rollback} 
et justifiez-les brièvement.  (1 point)

\begin{solution}
Réponse:  un \texttt{update} de \texttt{Spectacle} pour incrémenter \texttt{nbBilletsVendus},
et un \texttt{insert} dans \texttt{Billet}. Le \texttt{commit} vient à la fin car c'est seulement à ce moment
que la base est à nouveau dans un état cohérent. 
\end{solution}

\item Voici une exécution concurrente de deux réservations.
$$H=w_1(s)w_2(s)w_2(b_1)C_1w_1(b_2)C_2$$

    \item Montrez que cette exécution est sérialisable (1 point)
    
\begin{solution}
Réponse: un seul conflit
\end{solution}

    \item Quel est le résultat de l'application de l'algorithme de verrouillage à deux phases (1 point)?
    
\begin{solution}
Réponse: $T_1$ s'exécute avant $T_2$.
\end{solution}

\end{enumerate}

\end{document}

